%% ma: L12-B14-0005
%% lop: 12
%% chuong: 5
%% bai: 14
%% dang: 12.14.5
%% loai: TN4
%% muc: TH
%% dap_an: "B"
%% nguon: "(NBV)-ĐỀ SỐ 8-MỖI NGÀY 1 ĐỀ THI 2025.pdf (D08, MỖI NGÀY 1 ĐỀ - Đề số 8), Phần I Câu 11"
%% kien_thuc: "Bài 14 (phương trình mặt phẳng đi qua một điểm và có vectơ pháp tuyến), Bài 8 (tọa độ trung điểm, vectơ)"
%% trang_thai: chinh-thuc
%% ghi_chu: "Giáo viên đã duyệt 10/10/2026. Đáp án gốc B đúng."
\begin{ex}
Trong không gian $Oxyz$, cho hai điểm $A(-1;2;2)$ và $B(3;-2;-4)$. Phương trình mặt phẳng vuông góc và đi qua trung điểm của đoạn thẳng $AB$ là
\choice{$2x+2y-3z-5=0$}{\True $2x-2y-3z-5=0$}{$2x-2y+3z+1=0$}{$2x-2y-3z=0$}
\loigiai{Trung điểm của $AB$ là $M(1;0;-1)$. Mặt phẳng đi qua $M$ và vuông góc với $AB$ nhận $\overrightarrow{AB}=(4;-4;-6)$, hay $(2;-2;-3)$, làm vectơ pháp tuyến nên có phương trình
$2(x-1)-2(y-0)-3(z+1)=0$, tức là $2x-2y-3z-5=0$.}
\end{ex}
